\chapter{System Design and Implementation}

\section{System Architecture Overview}

FireWatch is designed as an integrated decision-support prototype for wildfire
monitoring in the Karabük regional context. The system combines weather-based
risk prediction, visual fire/smoke detection, persistent alert storage, and a
dashboard interface into one workflow. Its purpose is not to replace official
emergency systems, but to demonstrate how software, machine learning,
computer vision, and operator-facing tools can support earlier awareness and
more structured monitoring decisions.

The architecture is organised into several cooperating layers. The data and
model-input layer prepares daily weather inputs and engineered features. The
prediction layer applies the two-stage wildfire risk model, consisting of a
Stage 1 \texttt{HistGradientBoostingRegressor} and a Stage 2
\texttt{RandomForestClassifier}. The visual detection layer uses the active
YOLOv8 \texttt{best4.pt} detector to identify fire and smoke evidence from
camera sources. These layers are coordinated through a FastAPI backend, stored
through SQLite persistence, and presented through a React/TypeScript
dashboard. A drone-ready support layer provides mock, Tello-oriented, and demo
workflows for supervised prototype operation.

Prediction and detection are separated by design. The prediction module
estimates daily fire-weather risk from weather-derived inputs and produces
decision-support values such as predicted FWI and high-risk probability. The
detection module does not estimate FWI; instead, it examines camera frames and
confirms visible fire or smoke evidence. This separation keeps the numerical
risk model and the visual evidence pipeline understandable, testable, and
auditable.

The backend coordinates the system workflow and exposes the outputs needed by
the dashboard. Prediction runs, system state, and detection alerts are stored
so that operational information is not lost after a browser refresh or backend
restart. SQLite provides this persistence layer during the prototype stage by
preserving prediction history, the latest system state, and durable detection
alerts.

The dashboard acts as the operator-facing layer. It brings together risk
status, live monitoring views, detection alerts, model/system information, and
map context in a single interface. The drone-ready layer is included to show
how future camera streams and supervised patrol demos can connect to the same
monitoring workflow, but the current system does not claim full autonomous
production drone deployment.

\section{Backend and API Design}

The backend is implemented with FastAPI and functions as the coordination
layer for the FireWatch prototype. It connects the prediction pipeline,
monitoring services, detection alerting, scheduler status, drone-ready
support, and dashboard data responses. This design keeps the browser client
focused on presentation while the backend owns inference, persistence,
monitoring state, and API contracts.

The API is organised into high-level route groups. Risk and prediction
endpoints support manual risk checks and retrieval of the latest operational
prediction. Weather endpoints provide display-oriented weather information.
History endpoints expose stored prediction runs for audit and review. System
endpoints provide model metadata, health information, runtime configuration,
and scheduler status. Monitoring endpoints manage camera feeds, detection
notifications, and durable detection alerts. Drone endpoints expose
operator-controlled drone-ready status, stream, demo patrol, manual command,
and emergency-stop behaviours.

Manual and scheduled prediction requests use the same risk decision workflow
where possible. A risk check obtains model inputs, builds the required
features, runs the Stage 1 and Stage 2 models, applies the decision logic, and
persists the result. Scheduled runs are triggered by the backend scheduler,
while manual runs are initiated by an operator. Both run types are treated as
operational prediction events and can be reviewed later through the stored
history.

Camera streams and detection outputs are handled through the monitoring
services rather than through the prediction route group. Local cameras,
webcam/PC-camera feeds, and drone-ready video streams use the same general
fire/smoke detection path. When the detector raises an event, the alerting
service stores the label, confidence, source, timestamp, and optional snapshot
evidence for dashboard review.

The backend structure separates major responsibilities so that model
inference, camera monitoring, alert persistence, scheduler operation,
drone-ready controls, and response formatting remain manageable. This separation is
important for a prototype that combines several engineering areas, because it
reduces coupling between the risk model and the visual monitoring layer. The
design is local and prototype-oriented; it does not claim cloud deployment,
production-grade security, or integration with official emergency-response
infrastructure.

\section{Database and Persistence Design}

FireWatch uses SQLite as a lightweight persistence layer for the local
prototype. SQLite is suitable for this stage because it is serverless, easy to
inspect, and sufficient for a single-machine academic demonstration. The
active runtime database is stored at
\texttt{backend/outputs/karabuk\_fwi.db}.

The database preserves three main categories of information. The
\texttt{run\_history} table stores prediction runs, including run type,
timestamp, target date, predicted FWI, high-risk probability, final decision,
decision reason, thresholds, raw inputs, engineered features, run metadata,
and diagnostic details. The \texttt{system\_state} table stores small persistent state values,
such as the latest drone-ready policy state. The
\texttt{detection\_alerts} table stores durable fire/smoke alerts from the
monitoring layer.

Persistent storage is important because FireWatch is intended to support
operator review, not only live display. Prediction results should remain
available after refresh or restart so operators can inspect previous runs.
Detection alerts also need review status because an alert may be unread,
opened, marked as read, returned to unread, or hidden from normal dashboard
views. This makes persistence part of the decision-support design rather than
only a technical storage detail.

Detection alerts preserve the visual evidence metadata needed by the
dashboard. Stored fields include the fire/smoke label, confidence value,
source camera or device, timestamp, severity/message, detection count, detailed
detection payload, optional snapshot path, read/unread state, read timestamp,
soft-delete state, and deletion timestamp. Snapshot images are stored as files
and referenced from SQLite, while the database row keeps the operational
record available for the Detection Alerts interface.

The runtime database policy treats \texttt{backend/outputs/karabuk\_fwi.db}
as local runtime/demo data. It is ignored by Git and should not be committed,
because prediction runs, alert review state, soft-deleted alerts, and demo
activity change during normal use. Shared demonstration data should therefore
be prepared through seed scripts or clearly named snapshots rather than by
committing the active live database.

SQLite is an appropriate choice for a B.Sc. prototype and regional pilot, but
it is not presented as the final production database architecture. A
multi-user or deployed monitoring system would likely require a server-based
database such as PostgreSQL, stronger access control, backup policies, and
operational monitoring.

\section{Prediction and Detection Integration}

FireWatch integrates prediction and detection as two coordinated but logically
separate workflows. The prediction workflow estimates daily wildfire risk from
weather-derived features using the Stage 1 and Stage 2 models. The visual
detection workflow examines live camera frames using the YOLOv8
\texttt{best4.pt} detector and records visible fire or smoke evidence. Keeping
these workflows separate avoids mixing modelled risk with observed visual
evidence.

The Stage 1 prediction model produces a continuous predicted FWI value, while
the Stage 2 classifier supports the final high-risk decision near operationally
important thresholds. This output gives the operator a risk estimate for the
target date. The detection layer has a different role: it does not calculate
FWI or alter the prediction result, but checks whether camera sources contain
visible fire or smoke.

Both workflows meet in the backend and dashboard. Risk outputs are shown as
decision-support information, while detection alerts are shown as operational
evidence. This allows the operator to distinguish between a high-risk weather
situation and confirmed visual signs of fire or smoke. The design therefore
answers two different questions: whether conditions are dangerous, and whether
there is visible evidence that requires review.

When the prediction pipeline identifies high-risk conditions, the result can
support monitoring priority and drone-ready patrol recommendations. This does
not mean the system autonomously launches production drone missions. Instead,
the prediction result informs the operator and prepares the monitoring layer
for supervised or demo operation.

Detection alerts are persisted with practical review metadata, including the
camera or device source, fire/smoke label, confidence value, timestamp, and
snapshot or evidence reference when available. This gives the dashboard a
durable alert history and helps operators compare risk likelihood with visual
evidence from monitoring feeds.

\section{Frontend Dashboard Design}

The frontend dashboard is implemented with React and TypeScript. It is the
operator-facing layer of FireWatch and presents backend outputs in a structured
way. The dashboard is designed to support monitoring and review rather than to
duplicate backend logic. Prediction, persistence, detection, and scheduler
state remain owned by the backend.

The final dashboard structure is organised around approved operational views:
Overview / Live Monitoring, Detection Alerts, Risk Map, System and Model, and
System Flow. These views reduce the amount of information shown at once while
still giving the operator access to the main parts of the system.

The Overview / Live Monitoring view brings key information together. It is
intended to show current weather and risk context, the latest operational
status, live monitoring feeds, and the most important system summary. This
view acts as the starting point for daily use because it combines both risk
awareness and monitoring state.

The Detection Alerts view focuses on persistent fire/smoke evidence. It
supports review of alert labels, confidence values, source information,
timestamps, snapshots when available, read/unread state, and soft deletion.
This design makes alerts durable and reviewable rather than temporary
on-screen messages.

The Risk Map provides spatial context for the Karabük monitoring area and
supports planning-oriented interpretation of historical burned-area and risk
context. It is used as an operator awareness view rather than as a direct
replacement for the prediction model. The System and Model view summarizes
model, feature, threshold, and health information needed to understand the
current configuration. The System Flow view explains the end-to-end workflow
visually, showing how weather input, prediction, detection, persistence, and
dashboard review connect.

\section{Scheduler and Operational Run Design}

The scheduler supports repeated operational risk checks during the day. In the
current FireWatch design, scheduled checks are configured for 09:00, 11:00,
and 15:00 Istanbul time. Using the local operational timezone keeps scheduled
decisions aligned with the daily monitoring routine for Karabük.

Each scheduled run follows the same general data-to-decision path as an
operator-triggered manual run. The system obtains daily model inputs, builds
the required features, runs the Stage 1 and Stage 2 prediction workflow, and
applies the decision logic. This consistency helps avoid different behaviour
between automatic checks and manual checks.

Both manual and scheduled runs are persisted in SQLite. The
\texttt{run\_history} table stores prediction results and associated metadata,
allowing the dashboard to show the latest operational run and retain previous
decisions for later review. This shared persistence design also makes it
easier to compare manual and scheduled outputs for the same operational
period. This is important because risk decisions should be auditable rather
than disappearing after execution.

The scheduler design also prevents operational checks from being treated as a
temporary console-only process. Scheduled checks are registered as backend
jobs, exposed through system status, and connected to the same persistence path
used by manual checks. This makes the behaviour visible to the dashboard and
easier to verify during demonstrations.

The scheduled run design is still part of a local prototype. It supports
structured decision support and review, but it does not represent official
emergency deployment or certified operational wildfire command infrastructure.

\section{Drone-Ready Implementation}

FireWatch includes a drone-ready layer to prepare the prototype for supervised
aerial monitoring. This layer is designed to connect camera-based monitoring
with future drone streams while keeping physical drone operation under
operator control. It supports the project goal of being drone-ready without
claiming certified autonomous deployment.

The current implementation supports mock mode and Tello-oriented
configuration. Mock mode allows the drone workflow to be tested without real
hardware, which is useful for demonstrations, interface checks, and safe
development. Tello mode prepares the system for a configured drone adapter,
including connection status, video stream handling, battery/status reporting,
emergency stop support, and controlled demo behaviour when the required safety
flags are enabled.

Drone-ready support is connected to the monitoring and dashboard workflow.
The system can expose drone status, stream state, patrol recommendation state,
and demo patrol results to the operator interface. The same fire/smoke
detection path can process a drone-ready video source, so future aerial
monitoring can reuse the alerting and persistence design already used for
local camera feeds.

Physical movement is protected by safety gates. Manual controls are blocked by
default unless explicitly enabled. Takeoff or demo movement requires explicit
configuration and operator confirmation, and emergency stop support is
available through the drone-ready API layer. This design keeps the prototype
safe for academic demonstration and prevents risk prediction from directly
causing uncontrolled hardware movement.

High-risk prediction can support patrol recommendation or monitoring priority,
but the current project does not claim full autonomous production launch.
Full waypoint navigation, return-to-home behaviour, battery-aware mission
planning, regulatory compliance, and production safety validation remain
outside the current scope.

\section{Implementation Summary}

This chapter described how FireWatch is implemented as a working software
prototype. The system connects weather-based prediction, fire/smoke detection,
SQLite persistence, dashboard presentation, scheduled risk checks, and
drone-ready support into one decision-support workflow for the Karabük regional
pilot.

The FastAPI backend coordinates the main workflows and exposes the data needed
by the dashboard. SQLite preserves prediction history, system state, and
detection alerts so that operational information remains available for review.
The React/TypeScript dashboard presents risk information, alert evidence,
system/model status, workflow context, and spatial awareness to the operator.

The implementation is modular: prediction, detection, persistence, scheduler
operation, drone-ready support, and dashboard views are separated into clear
responsibilities. This makes the prototype easier to test, explain, and
extend. It also supports the methodology described in Chapter 3 and prepares
the basis for the results and discussion presented in Chapter 5.
